\documentclass[tikz,border=3pt]{standalone}
\usepackage[utf8]{inputenc}
\usepackage[T1]{fontenc}
\usepackage{lmodern}
\renewcommand{\familydefault}{\sfdefault}
\usepackage{sansmath}\sansmath
\usepackage{chemfig}
\usepackage{pgfplots}\pgfplotsset{compat=1.18,/pgf/number format/use comma,/pgf/number format/1000 sep={.}}
\usetikzlibrary{arrows.meta,calc,positioning,decorations.pathmorphing,patterns}
% paleta del sitio
\definecolor{nar}{HTML}{E8680C}\definecolor{narc}{HTML}{FFB35C}\definecolor{hielo}{HTML}{FFF3E8}
\definecolor{tinta}{HTML}{1D1D1F}\definecolor{gris}{HTML}{6E6E73}\definecolor{lila}{HTML}{7C5CF0}
\definecolor{azul}{HTML}{2F6FD6}\definecolor{menta}{HTML}{17A673}\definecolor{coral}{HTML}{E8342A}
\usepackage[european]{circuitikz}
\begin{document}
\begin{circuitikz}[font=\footnotesize]
\draw[thick] (0,0) to[battery1, l=pila] (0,3) to[rmeter, t=A, l=amperímetro, color=azul] (4,3) to[lamp, l_=ampolleta] (4,0) -- (0,0);
\draw[thick] (4,3) -- (6.6,3) to[rmeter, t=V, l=voltímetro, color=menta!70!black] (6.6,0) -- (4,0);
\node[align=center,text=azul,font=\scriptsize] at (2,4.0) {en serie: mide la intensidad de corriente};
\node[align=center,text=menta!70!black,font=\scriptsize] at (5.3,-0.7) {en paralelo: mide la diferencia de potencial de la ampolleta};
\end{circuitikz}
\end{document}
